\begin{proposition}[\texorpdfstring{$\dot H^{1/2}(\RR)$}{Hdot1/2(R)} 能量的闭式 Gram 核与端点缺陷势]\label{prop:xi-endpoint-jensen-defect-h12-gram-kernel}
\leanverified{paper\_xi\_endpoint\_jensen\_defect\_h12\_gram\_kernel}
取有限参数族 \(\{(\gamma_j,\delta_j,m_j)\}_{j=1}^{J}\subset \RR\times(0,1)\times\RR\)，并令一维 Poisson 核为
\[
\mathsf P_s(x):=\frac{1}{\pi}\frac{s}{x^2+s^2}\qquad(s>0).
\]
定义端点缺陷型势函数
\[
\Psi_{(\gamma,\delta)}(x)
:=\pi\int_{1-\delta}^{1+\delta}\mathsf P_s(x-\gamma)\,\dd s
=\frac12\log\frac{(x-\gamma)^2+(1+\delta)^2}{(x-\gamma)^2+(1-\delta)^2},
\]
并令
\[
\mathcal D(x):=\sum_{j=1}^{J} m_j\,\Psi_{(\gamma_j,\delta_j)}(x).
\]
则 \(\mathcal D\in \dot H^{1/2}(\RR)\)，且其齐次半范数满足完全显式的双重求和公式
\begin{equation}\label{eq:xi-endpoint-jensen-defect-h12-gram-kernel}
\boxed{\;
\|\mathcal D\|_{\dot H^{1/2}}^2
=\frac{\pi}{2}\sum_{j,k=1}^{J} m_jm_k\,\mathbf K\!\bigl((\gamma_j,\delta_j),(\gamma_k,\delta_k)\bigr)\;}
\end{equation}
其中 Gram 核为
\[
\mathbf K\!\bigl((\gamma,\delta),(\gamma',\delta')\bigr)
:=\log\!\Bigg[
\frac{\bigl((2+\delta'-\delta)^2+(\gamma-\gamma')^2\bigr)\bigl((2+\delta-\delta')^2+(\gamma-\gamma')^2\bigr)}
{\bigl((2-\delta-\delta')^2+(\gamma-\gamma')^2\bigr)\bigl((2+\delta+\delta')^2+(\gamma-\gamma')^2\bigr)}
\Bigg].
\]
\end{proposition}
\begin{proof}
取傅里叶约定 \(\widehat f(\xi)=\int_{\RR}f(x)e^{-i\xi x}\,\dd x\)。
由 \(\widehat{\mathsf P_s(\cdot-\gamma)}(\xi)=e^{-s|\xi|}e^{-i\gamma\xi}\) 得
\[
\widehat{\Psi_{(\gamma,\delta)}}(\xi)
=\pi e^{-i\gamma\xi}\int_{1-\delta}^{1+\delta}e^{-s|\xi|}\,\dd s
=\pi e^{-i\gamma\xi}\frac{e^{-(1-\delta)|\xi|}-e^{-(1+\delta)|\xi|}}{|\xi|}.
\]
故
\[
\widehat{\mathcal D}(\xi)
=\pi\sum_{j=1}^{J} m_j e^{-i\gamma_j\xi}\frac{e^{-(1-\delta_j)|\xi|}-e^{-(1+\delta_j)|\xi|}}{|\xi|}.
\]
再用齐次半范数表示
\[
\|\mathcal D\|_{\dot H^{1/2}}^2=\frac{1}{2\pi}\int_{\RR}|\xi|\,|\widehat{\mathcal D}(\xi)|^2\,\dd\xi,
\]
展开并利用偶对称化得到
\[
\|\mathcal D\|_{\dot H^{1/2}}^2
=\pi\sum_{j,k=1}^{J}m_jm_k\int_0^\infty \frac{\cos((\gamma_j-\gamma_k)\xi)}{\xi}
\bigl(e^{-(1-\delta_j)\xi}-e^{-(1+\delta_j)\xi}\bigr)\bigl(e^{-(1-\delta_k)\xi}-e^{-(1+\delta_k)\xi}\bigr)\,\dd\xi.
\]
对任意 \(0<a<b\) 与 \(c\in\RR\) 有标准恒等式
\[
\int_0^\infty \frac{e^{-a\xi}-e^{-b\xi}}{\xi}\cos(c\xi)\,\dd\xi
=\frac12\log\frac{b^2+c^2}{a^2+c^2},
\]
将乘积差分写为四项差分之线性组合并逐项套用该恒等式，即得 \eqref{eq:xi-endpoint-jensen-defect-h12-gram-kernel} 的闭式核表达。证毕。
\end{proof}

\begin{proposition}[\texorpdfstring{$\dot H^{1/2}$}{Hdot1/2} 刚性：核的严格正定性判别]\label{prop:xi-endpoint-jensen-defect-h12-strict-pd}
\leanverified{paper\_xi\_endpoint\_jensen\_defect\_h12\_strict\_pd}
沿用命题 \ref{prop:xi-endpoint-jensen-defect-h12-gram-kernel} 的记号。
若参数对 \((\gamma_j,\delta_j)\) 两两不同，则 Gram 矩阵
\(\mathbf K_{jk}:=\mathbf K\!\bigl((\gamma_j,\delta_j),(\gamma_k,\delta_k)\bigr)\)
严格正定；从而对任何非零实系数向量 \((m_j)\neq 0\) 有
\[
\|\mathcal D\|_{\dot H^{1/2}}^2>0.
\]
特别地，若 \(m_j\ge 0\) 且 \(\mathcal D(x)\equiv 0\)（几乎处处），则必有 \(m_j=0\) 全部成立。
\end{proposition}
\begin{proof}
由 \(\dot H^{1/2}\) 的 Hilbert 结构，\(\|\mathcal D\|_{\dot H^{1/2}}^2\ge 0\)。
若 \(\|\mathcal D\|_{\dot H^{1/2}}^2=0\)，则 \( |\xi|^{1/2}\widehat{\mathcal D}(\xi)=0 \) 于几乎处处，从而 \(\widehat{\mathcal D}(\xi)=0\) 于 \(\xi\neq 0\)。
由命题 \ref{prop:xi-endpoint-jensen-defect-h12-gram-kernel} 的傅里叶表达式，得到对所有 \(\xi>0\),
\[
\sum_{j=1}^{J} m_j\,e^{-i\gamma_j\xi}\bigl(e^{-(1-\delta_j)\xi}-e^{-(1+\delta_j)\xi}\bigr)=0.
\]
将左端视为右半平面 \(\{z\in\CC:\Re z>0\}\) 上的解析函数
\[
F(z):=\sum_{j=1}^{J}m_j\bigl(e^{-((1-\delta_j)+i\gamma_j)z}-e^{-((1+\delta_j)+i\gamma_j)z}\bigr),
\]
其在 \((0,\infty)\) 上恒为零，故由恒等定理知 \(F\equiv 0\)。
由于指数率 \(((1\pm\delta_j)+i\gamma_j)\) 在 \((\gamma_j,\delta_j)\) 两两不同条件下两两不同，有限指数和的线性无关性推出所有系数均为零，即 \(m_j=0\)（\(1\le j\le J\)）。证毕。
\end{proof}

\begin{proposition}[交比表示与 \texorpdfstring{$PSL(2,\RR)$}{PSL(2,R)} Möbius 不变性]\label{prop:xi-endpoint-jensen-defect-h12-crossratio-mobius}
\leanverified{paper\_xi\_endpoint\_jensen\_defect\_h12\_crossratio\_mobius}
在命题 \ref{prop:xi-endpoint-jensen-defect-h12-gram-kernel} 的设定下，令上半平面点
\[
z_j^\pm:=\gamma_j+i(1\pm\delta_j)\in\mathbb H.
\]
则核 \(\mathbf K_{jk}\) 可写为交比模长：
\[
\mathbf K_{jk}
=2\log\left|
\frac{(z_j^- - \overline{z_k^+})(z_j^+ - \overline{z_k^-})}{(z_j^- - \overline{z_k^-})(z_j^+ - \overline{z_k^+})}
\right|.
\]
因此 \(\|\mathcal D\|_{\dot H^{1/2}}^2\) 具有等价表示
\[
\|\mathcal D\|_{\dot H^{1/2}}^2
=\pi\sum_{j,k=1}^{J} m_jm_k\,
\log\left|
\frac{(z_j^- - \overline{z_k^+})(z_j^+ - \overline{z_k^-})}{(z_j^- - \overline{z_k^-})(z_j^+ - \overline{z_k^+})}
\right|.
\]
并且对任意 \(\phi\in PSL(2,\RR)\) 满足
\[
\log\left|
\frac{(\phi(z_j^-)-\overline{\phi(z_k^+)})(\phi(z_j^+)-\overline{\phi(z_k^- )})}{(\phi(z_j^-)-\overline{\phi(z_k^-)})(\phi(z_j^+)-\overline{\phi(z_k^+)})}
\right|
=
\log\left|
\frac{(z_j^- - \overline{z_k^+})(z_j^+ - \overline{z_k^-})}{(z_j^- - \overline{z_k^-})(z_j^+ - \overline{z_k^+})}
\right|.
\]
\end{proposition}
\begin{proof}
注意到
\[
|z_j^\pm-\overline{z_k^\mp}|^2
=(\gamma_j-\gamma_k)^2+\bigl((1\pm\delta_j)+(1\mp\delta_k)\bigr)^2
=(\gamma_j-\gamma_k)^2+(2\pm\delta_j\mp\delta_k)^2,
\]
以及
\(
|z_j^\pm-\overline{z_k^\pm}|^2=(\gamma_j-\gamma_k)^2+(2\pm\delta_j\pm\delta_k)^2
\)。
将其代回命题 \ref{prop:xi-endpoint-jensen-defect-h12-gram-kernel} 的闭式核即可得到交比表达式。
\par
对 \(\phi(z)=\frac{az+b}{cz+d}\)（\(a,b,c,d\in\RR\)、\(ad-bc>0\)），有 \(\phi(\overline z)=\overline{\phi(z)}\)，并且交比在 Möbius 变换下不变。
将
\(\alpha=z_j^-\)、\(\beta=\overline{z_k^+}\)、\(\gamma=z_j^+\)、\(\delta=\overline{z_k^-}\)
代入并取模长与对数，得到所述不变性。证毕。
\end{proof}

\begin{theorem}[Poisson 平移下的 Loewner 单调 Gram 核与耗散导数]\label{thm:xi-endpoint-psi-poisson-shift-loewner-monotone}
\leanverified{paper\_xi\_endpoint\_psi\_poisson\_shift\_loewner\_monotone}
沿用命题 \ref{prop:xi-endpoint-jensen-defect-h12-gram-kernel} 的记号，并对每个 \(t\ge 0\) 定义 Poisson 平移势
\[
\Psi_{(\gamma,\delta)}^{(t)}:=\mathsf P_t*\Psi_{(\gamma,\delta)}.
\]
对任意有限点集 \(\{(\gamma_j,\delta_j)\}_{j=1}^{J}\)，定义 \(\dot H^{1/2}\) 的 Gram 矩阵族
\[
K_t(j,k):=\frac{2}{\pi}\,\bigl\langle \Psi_{(\gamma_j,\delta_j)}^{(t)},\,\Psi_{(\gamma_k,\delta_k)}^{(t)}\bigr\rangle_{\dot H^{1/2}(\RR)}
\qquad(1\le j,k\le J).
\]
则：
\begin{enumerate}
  \item 对任意 \(0\le t_1<t_2\) 有 Loewner 单调性 \(K_{t_1}-K_{t_2}\succeq 0\)。
  \item 对任意 \(m\in\RR^{J}\)，函数 \(t\mapsto m^{\mathsf T}K_t m\) 在 \((0,\infty)\) 上可微且满足完全耗散导数公式
  \[
  \boxed{\;
  \frac{\dd}{\dd t}\bigl(m^{\mathsf T}K_t m\bigr)
  =-\frac{2}{\pi^2}\int_{\RR}|\xi|^2\,
  \left|\sum_{j=1}^{J}m_j\,\widehat{\Psi_{(\gamma_j,\delta_j)}}(\xi)\right|^2
  e^{-2t|\xi|}\,\dd\xi\ \le\ 0. \;}
  \]
\end{enumerate}
\end{theorem}
\begin{proof}
由 \(\widehat{\mathsf P_t*f}=e^{-t|\xi|}\widehat f\) 以及命题 \ref{prop:xi-endpoint-jensen-defect-h12-gram-kernel} 中的 Fourier 表达式，
对任意 \(m\in\RR^J\) 令 \(\mathcal D:=\sum_{j=1}^{J}m_j\Psi_{(\gamma_j,\delta_j)}\)，则
\[
m^{\mathsf T}K_t m
=\frac{2}{\pi}\,\|\mathsf P_t*\mathcal D\|_{\dot H^{1/2}}^2
=\frac{2}{\pi}\cdot\frac{1}{2\pi}\int_{\RR}|\xi|\,|\widehat{\mathcal D}(\xi)|^2\,e^{-2t|\xi|}\,\dd\xi
=\frac{1}{\pi^2}\int_{\RR}|\xi|\,|\widehat{\mathcal D}(\xi)|^2\,e^{-2t|\xi|}\,\dd\xi,
\]
其中 \(\widehat{\mathcal D}(\xi)=\sum_j m_j\widehat{\Psi_{(\gamma_j,\delta_j)}}(\xi)\)。
于是当 \(t_1<t_2\) 时，
\[
m^{\mathsf T}(K_{t_1}-K_{t_2})m
=\frac{1}{\pi^2}\int_{\RR}|\xi|\,|\widehat{\mathcal D}(\xi)|^2\bigl(e^{-2t_1|\xi|}-e^{-2t_2|\xi|}\bigr)\,\dd\xi\ \ge\ 0,
\]
对一切 \(m\) 成立，故 \(K_{t_1}-K_{t_2}\succeq 0\)。
对导数，指数权重给出支配收敛，故可将微分移入积分并得
\[
\frac{\dd}{\dd t}\bigl(m^{\mathsf T}K_t m\bigr)
=-\frac{2}{\pi^2}\int_{\RR}|\xi|^2\,|\widehat{\mathcal D}(\xi)|^2\,e^{-2t|\xi|}\,\dd\xi,
\]
将 \(\widehat{\mathcal D}\) 展开即得所示公式（并因被积项非负而导数非正）。证毕。
\end{proof}

\begin{theorem}[Poisson 平移下的单极通道锐渐近]\label{thm:xi-endpoint-psi-poisson-monopole-asymptotic}
\leanverified{paper\_xi\_endpoint\_psi\_poisson\_monopole\_asymptotic}
沿用命题 \ref{prop:xi-endpoint-jensen-defect-h12-gram-kernel} 的记号，令
\(
\mathcal D=\sum_{j=1}^{J}m_j\Psi_{(\gamma_j,\delta_j)}
\)
并记 \(\mathcal D_t:=\mathsf P_t*\mathcal D\)。
令
\[
Q_0:=\sum_{j=1}^{J}m_j\delta_j.
\]
若 \(Q_0\neq 0\)，则极限存在且满足
\[
\boxed{\;
\lim_{t\to\infty}t^{2}\,\|\mathcal D_t\|_{\dot H^{1/2}(\RR)}^{2}
=\pi\,Q_0^{2}. \;}
\]
\end{theorem}
\begin{proof}
由 Fourier 表达式
\(
\|\mathcal D_t\|_{\dot H^{1/2}}^{2}=\frac{1}{2\pi}\int_{\RR}|\xi|\,|\widehat{\mathcal D}(\xi)|^{2}e^{-2t|\xi|}\,\dd\xi
\)。
另一方面，
\(
\widehat{\Psi_{(\gamma,\delta)}}(\xi)=\pi e^{-i\gamma\xi}\frac{e^{-(1-\delta)|\xi|}-e^{-(1+\delta)|\xi|}}{|\xi|}
\)
给出
\(
\widehat{\Psi_{(\gamma,\delta)}}(0)=2\pi\delta
\)，从而 \(\widehat{\mathcal D}(0)=2\pi Q_0\)。
作变量代换 \(\xi=y/t\)，得
\[
t^2\|\mathcal D_t\|_{\dot H^{1/2}}^{2}
=\frac{1}{2\pi}\int_{\RR}|y|\,|\widehat{\mathcal D}(y/t)|^{2}e^{-2|y|}\,\dd y.
\]
由 \(\widehat{\mathcal D}\) 在 \(0\) 处连续且 \(e^{-2|y|}\) 支配收敛，令 \(t\to\infty\) 得
\[
\lim_{t\to\infty}t^2\|\mathcal D_t\|_{\dot H^{1/2}}^{2}
=\frac{1}{2\pi}\,|\widehat{\mathcal D}(0)|^{2}\int_{\RR}|y|e^{-2|y|}\,\dd y
=\frac{1}{2\pi}\,(2\pi Q_0)^{2}\cdot \frac12
=\pi Q_0^2.
\]
证毕。
\end{proof}

\begin{theorem}[Poisson 平移下的偶极通道与 \texorpdfstring{$t^{-4}$}{t^-4} 加速]\label{thm:xi-endpoint-psi-poisson-dipole-asymptotic}
沿用定理 \ref{thm:xi-endpoint-psi-poisson-monopole-asymptotic} 的记号。
设
\[
Q_0:=\sum_{j=1}^{J}m_j\delta_j=0,
\qquad
Q_1:=\sum_{j=1}^{J}m_j\delta_j\gamma_j\neq 0.
\]
则极限存在且满足
\[
\boxed{\;
\lim_{t\to\infty}t^{4}\,\|\mathcal D_t\|_{\dot H^{1/2}(\RR)}^{2}
=\frac{3\pi}{2}\,Q_1^{2}. \;}
\]
并且若进一步 \(Q_0=Q_1=0\)，则有严格更强衰减 \(\|\mathcal D_t\|_{\dot H^{1/2}}^{2}=o(t^{-4})\)。
\leanverified{paper\_xi\_endpoint\_psi\_poisson\_dipole\_asymptotic}
\end{theorem}
\begin{proof}
令
\[
A_\delta(u):=\frac{e^{-(1-\delta)u}-e^{-(1+\delta)u}}{u}\qquad(u>0).
\]
则
\(
\widehat{\Psi_{(\gamma,\delta)}}(\xi)=\pi e^{-i\gamma\xi}A_\delta(|\xi|)
\)。
当 \(u\downarrow 0\) 时有展开 \(A_\delta(u)=2\delta-2\delta u+O(u^2)\)。
因此
\[
\widehat{\mathcal D}(\xi)
=\pi\sum_{j=1}^{J}m_j e^{-i\gamma_j\xi}\bigl(2\delta_j-2\delta_j|\xi|+O(|\xi|^2)\bigr).
\]
若 \(Q_0=\sum_j m_j\delta_j=0\)，则常数项与 \(|\xi|\) 项同时消去，且用 \(e^{-i\gamma_j\xi}=1-i\gamma_j\xi+O(\xi^2)\) 得
\[
\widehat{\mathcal D}(\xi)=-2\pi i\xi\,Q_1+O(|\xi|^2),
\qquad
|\widehat{\mathcal D}(\xi)|^2=4\pi^2\xi^2Q_1^2+O(|\xi|^3).
\]
代入
\(
\|\mathcal D_t\|_{\dot H^{1/2}}^{2}=\frac{1}{2\pi}\int_{\RR}|\xi|\,|\widehat{\mathcal D}(\xi)|^{2}e^{-2t|\xi|}\,\dd\xi
\)，作缩放 \(\xi=y/t\) 并用支配收敛，得到
\[
\lim_{t\to\infty}t^4\|\mathcal D_t\|_{\dot H^{1/2}}^{2}
=\frac{1}{2\pi}\cdot 4\pi^2Q_1^2\int_{\RR}|y|^3e^{-2|y|}\,\dd y
=\frac{1}{2\pi}\cdot 4\pi^2Q_1^2\cdot \frac34
=\frac{3\pi}{2}Q_1^2.
\]
若进一步 \(Q_1=0\)，则 \(\widehat{\mathcal D}(\xi)=O(|\xi|^2)\)，同样缩放给出 \(\|\mathcal D_t\|_{\dot H^{1/2}}^{2}=o(t^{-4})\)。证毕。
\end{proof}

\begin{corollary}[矩阵值完全单调性]\label{cor:xi-endpoint-psi-poisson-complete-monotonicity}
在定理 \ref{thm:xi-endpoint-psi-poisson-shift-loewner-monotone} 的设定下，
矩阵函数 \(t\mapsto K_t\) 在 \((0,\infty)\) 上 \(C^\infty\)，并且对所有整数 \(n\ge 0\) 都有
\[
\boxed{\;
(-1)^n\frac{\dd^n}{\dd t^n}K_t\ \succeq\ 0\qquad(\text{Loewner 序}).\;}
\]
等价地，对任意 \(m\in\RR^J\)，标量函数 \(t\mapsto m^{\mathsf T}K_t m\) 为完全单调函数，因此存在唯一有限正测度 \(\mu_m\) 使
\[
m^{\mathsf T}K_t m=\int_{0}^{\infty}e^{-2ts}\,\dd\mu_m(s).
\]
\end{corollary}
\begin{proof}
由定理 \ref{thm:xi-endpoint-psi-poisson-shift-loewner-monotone} 的积分表示，
\[
m^{\mathsf T}K_t m=\frac{1}{\pi^2}\int_{\RR}|\xi|\,|\widehat{\mathcal D}(\xi)|^2e^{-2t|\xi|}\,\dd\xi,
\qquad \mathcal D=\sum_j m_j\Psi_{(\gamma_j,\delta_j)}.
\]
对 \(t\) 取 \(n\) 阶导数并将导数移入积分，得到
\[
(-1)^n\frac{\dd^n}{\dd t^n}\bigl(m^{\mathsf T}K_t m\bigr)
=\frac{2^n}{\pi^2}\int_{\RR}|\xi|^{n+1}\,|\widehat{\mathcal D}(\xi)|^2e^{-2t|\xi|}\,\dd\xi\ \ge\ 0,
\]
对一切 \(m\) 成立，等价于 \((-1)^n\dd^nK_t/\dd t^n\succeq 0\)。
测度表示由 Bernstein 定理应用于完全单调函数 \(t\mapsto m^{\mathsf T}K_t m\) 得到。证毕。
\end{proof}

\endinput
